% Compile: xelatex Lecture_7_Level_Pressure_Math.tex (twice)
\documentclass[aspectratio=169,11pt]{beamer}
\usepackage{fontspec}
\usepackage{polyglossia}
\setdefaultlanguage{ukrainian}
\setmainfont{DejaVu Sans}
\setsansfont{DejaVu Sans}
\newfontfamily\cyrillicfont{DejaVu Sans}
\newfontfamily\cyrillicfontsf{DejaVu Sans}
\newfontfamily\cyrillicfonttt{DejaVu Sans Mono}
\setmonofont{DejaVu Sans Mono}
\usepackage{amsmath,amssymb,booktabs}
\definecolor{ink}{HTML}{183341}
\definecolor{accent}{HTML}{087C83}
\setbeamercolor{normal text}{fg=ink,bg=white}
\setbeamercolor{frametitle}{fg=accent}
\setbeamertemplate{navigation symbols}{}
\setbeamertemplate{footline}[frame number]
\title{Лекція 7. Геометричні розміри, рівень і тиск}
\subtitle{Математичний конспект і практичні завдання}
\date{Новий план лекцій 7--14, редакція 3}
\begin{document}
\begin{frame}\titlepage\end{frame}
\begin{frame}{Умови та позначення}
Цей конспект супроводжує презентацію лекції 7. Числові приклади й паспорти задач синтетичні.
\begin{itemize}
\item $h,H,d$ --- висоти та відстані, м; $p$ --- тиск, Па.
\item $\rho$ --- густина, кг/м$^3$; $g=9{,}81$ м/с$^2$ у навчальних прикладах.
\item $1$ бар $=10^5$ Па; $1$ кПа $=10^3$ Па; $1$ мА $=10^{-3}$ А.
\item Геометрична база, опорний тиск і одиниці є частиною вхідних даних.
\end{itemize}
Гранична похибка, стандартна невизначеність і роздільна здатність мають різний зміст. Тут клас використовується лише за явно заданим правилом.
\end{frame}
\section{Геометричні розміри та механічне вимірювання рівня}
\begin{frame}{Геометричні розміри та механічне вимірювання рівня}
За спільної вертикальної бази $h=H-d$, де $H$ --- висота сенсора, $d$ --- відстань до поверхні. Для нахиленого контакту вимірювана довжина потребує геометричної проєкції.
\[
 F_A=\rho gV_{\mathrm{зан}},\qquad [F_A]=\mathrm{N}.
\]
У вільного поплавка без додаткових сил $F_A=mg$. Зменшення густини збільшує потрібний занурений об'єм: $V_{\mathrm{зан}}=m/\rho$.
\[
 T=mg-\rho gV_{\mathrm{зан}},\qquad \Delta T=-\rho gA\Delta h.
\]
Остання формула діє для вертикального буйка зі сталим перерізом $A$ (м$^2$) на ділянці часткового занурення. Тертя, налипання та сила приводу потребують додаткових членів.
\end{frame}
\begin{frame}{Визначення: лінійна характеристика рівня}
Для двох калібрувальних точок $(I_0,h_0)$ і $(I_1,h_1)$:
\[
 h=h_0+\frac{I-I_0}{I_1-I_0}(h_1-h_0),\qquad I_1\ne I_0.
\]
У задачі $h_0=0$ м, $h_1=2$ м, $I_0=4$ мА, $I_1=20$ мА:
\[
 h=0{,}125\;\frac{\mathrm{м}}{\mathrm{мА}}(I-4\;\mathrm{мА}).
\]
Перевірка кінців: $h(4\;\mathrm{мА})=0$ м, $h(20\;\mathrm{мА})=2$ м.
\[
 I=4\;\mathrm{мА}+8\;\frac{\mathrm{мА}}{\mathrm{м}}h.
\]
Підстановка прямого перетворення в обернене повертає $I$. Екстрапольоване число поза інтервалом не є підтвердженим результатом калібрування.
\end{frame}
\section{Гідростатичні та безконтактні рівнеміри}
\begin{frame}{Гідростатичні та безконтактні рівнеміри}
Для нерухомої однорідної рідини, без поправок імпульсних ліній:
\[
 p_{\mathrm{низ}}=p_{\mathrm{газ}}+\rho gh,\qquad h=\frac{p_{\mathrm{низ}}-p_{\mathrm{газ}}}{\rho g}.
\]
Відкритий бак: $p_{\mathrm{газ}}=p_{\mathrm{атм}}$. Закритий бак: тиск газу треба вимірювати або підводити до опорної камери дифперетворювача.
\[
 \frac{\mathrm{Pa}}{(\mathrm{kg/m^3})(\mathrm{m/s^2})}=\mathrm{m}.
\]
Якщо густина залежить від глибини, $\Delta p=g\int_0^h\rho(z)\,dz$; проста формула потребує сталої або обґрунтованої середньої густини.
\end{frame}
\begin{frame}{Розрахунок і перевірка рівня}
\[
 h=\frac{19620\;\mathrm{Pa}}{1000\;\mathrm{kg/m^3}\cdot9{,}81\;\mathrm{m/s^2}}=2{,}00\;\mathrm{м}.
\]
За $\rho=800$ кг/м$^3$ той самий перепад дає $h=2{,}50$ м. Якщо використати густину води замість фактичної, оцінка буде занижена на $0{,}50$ м.

Зворотна перевірка: $\Delta p=\rho gh=19620$ Па. Для бака висотою $H=3$ м і води допустимий геометричний інтервал відповідає $0\le\Delta p\le29430$ Па.

У закритому баку за $p_{\mathrm{газ,g}}=30000$ Па та $h=2$ м нижній надлишковий тиск становить $49620$ Па. Віднімати треба тиски з однаковою опорою.
\end{frame}
\begin{frame}{Визначення: час прольоту й FMCW}
\[
 d=\frac{vt}{2},\qquad h=H-\frac{vt}{2},\qquad t=\frac{2(H-h)}v.
\]
$v$ --- швидкість хвилі (м/с), $t$ --- час подвійного шляху (с). У задачі ультразвук: $v=343$ м/с. Радар у повітрі: $v\approx3\cdot10^8$ м/с.

Для $H=3$ м, $h=2$ м: $t\approx0{,}005831$ с для ультразвуку, близько $6{,}67$ нс для радара. Підміна швидкості змінює результат на багато порядків.
\[
 f_b\approx Kt,\qquad d\approx\frac{vf_b}{2K},\quad [K]=\mathrm{Hz/s},\quad [f_b]=\mathrm{Hz}.
\]
FMCW: нерухома поверхня, лінійна частотна розгортка, коректно виділене відлуння. Мертва зона $d_{\min}$ обмежує рівень: $h\le H-d_{\min}$.
\end{frame}
\begin{frame}{Причини розбіжності двох оцінок}
\begin{tabular}{p{.30\textwidth}p{.58\textwidth}}\toprule
Канал & Що перевіряємо \\\midrule
Гідростатичний & Густину, опорний газовий тиск, висоти підключень, заповнення ліній \\
Продувний & Положення виходу, витрату газу, втрати, закупорювання \\
Ультразвуковий & Швидкість у газі, температуру, мертву зону, піну \\
Радарний & Вибір відлуння, геометрію, антену, властивості середовища \\\bottomrule
\end{tabular}
\medskip
$r=h_p-h_t$ (м) --- розбіжність, а не автоматично похибка одного з каналів. Просте усереднення не усуває невідому систематичну причину.
\end{frame}
\section{Рідинні та деформаційні манометри}
\begin{frame}{Рідинні та деформаційні манометри}
\[
 p_{\mathrm{abs}}=p_{\mathrm{g}}+p_{\mathrm{atm}},\qquad \Delta p=p_1-p_2.
\]
Усі тиски в Па. Від'ємний $p_{\mathrm{g}}$ означає тиск нижче місцевого атмосферного, а не від'ємний абсолютний тиск.
\[
 \Delta p=\rho_mg\Delta h.
\]
Спрощення для газу над рідиною, густиною газу нехтуємо. $\Delta h$ --- вертикальна різниця двох поверхонь, м; $\rho_m$ --- густина манометричної рідини, кг/м$^3$.

Для відповідної схеми з іншою рідиною над обома колінами:
$\Delta p=(\rho_m-\rho_f)g\Delta h$. Для похилої трубки вертикальна проєкція довжини $l$: $\Delta h=l\sin\theta$.
\end{frame}
\begin{frame}{Визначення: пружне перетворення та гістерезис}
У лінійному навчальному наближенні мембрана:
\[
 F=\Delta p A_{\mathrm{eff}},\qquad F=k_s x,\qquad
 x=\frac{A_{\mathrm{eff}}}{k_s}\Delta p.
\]
$F$ у Н, $A_{\mathrm{eff}}$ у м$^2$, $k_s$ у Н/м, $x$ у м. Реальні мембрани можуть мати змінну ефективну площу та нелінійну жорсткість.

Трубка Бурдона використовує зміну форми трубки овального перерізу; сильфон --- осьову деформацію гофрованої оболонки. У силовому зрівноваженні $F_{\mathrm{комп}}\approx\Delta p A_{\mathrm{eff}}$.

Гістерезис у контрольній точці характеризує різниця $I_{\uparrow}(p)-I_{\downarrow}(p)$ у спільних одиницях. Перевірка тільки прямого ходу його не визначає.
\end{frame}
\begin{frame}{Визначення: межа за класом у навчальній задачі}
Прийняте правило:
\[
 B=\frac{k}{100}(p_{\max}-p_{\min}).
\]
$k=1{,}5$ --- числове значення відсотка; $B$ та межі шкали в барах.
\begin{center}\begin{tabular}{ccc}\toprule
Шкала, бар & $B$, бар & За вимоги $B\le0{,}07$ бар \\\midrule
0--4 & 0,06 & проходить \\
0--6 & 0,09 & не проходить \\
0--10 & 0,15 & не проходить \\\bottomrule
\end{tabular}\end{center}
Тиск $2{,}5$ бара входить у всі три шкали. Придатність за температурою, пульсаціями та середовищем цим розрахунком не встановлено.
\end{frame}
\section{Електричні й інтелектуальні перетворювачі тиску}
\begin{frame}{Електричні й інтелектуальні перетворювачі тиску}
\[
 \frac{\Delta R}{R}=k\varepsilon,\qquad \varepsilon=\frac{\Delta l}{l}.
\]
$R,\Delta R$ в Ом; $k,\varepsilon$ безрозмірні. Міст перетворює зміни опорів на напругу. Температурну компенсацію перевіряють за характеристиками всього каналу.
\[
 C\approx\frac{\varepsilon_0\varepsilon_r A}{d}.
\]
$C$ у Ф; $\varepsilon_0$ у Ф/м; $\varepsilon_r$ безрозмірна; $A$ у м$^2$; $d$ у м. При $d=d_0-x$ і $|x|\ll d_0$: $C\approx C_0(1+x/d_0)$.

Індуктивне зчитування використовує переміщення осердка та зміну індуктивності (Гн) або взаємної індуктивності. Тип сенсора сам не визначає час відгуку повного каналу.
\end{frame}
\begin{frame}{Визначення: заряд і часові обмеження}
Спрощена п'єзоелектрична модель: $Q=s_FF$, де $Q$ у Кл, $s_F$ у Кл/Н, $F$ у Н.
\[
 U(t)=U_0 e^{-t/\tau},\qquad [U]=\mathrm{V},\quad [t]=[\tau]=\mathrm{s}.
\]
Спад за сталого навантаження ілюструє витік заряду. Придатність для квазістатичного вимірювання залежить від тривалості, дрейфу, підсилювача та потрібної межі.

Для каналу зі сталим кінцевим відгуком час усталення задають строго:
\[
 t_s=\inf\{t_0\ge0:\ |y(t)-y_\infty|\le\delta\ \text{для всіх }t\ge t_0\}.
\]
$\delta$ у тих самих одиницях, що $y$. Перше входження у смугу не гарантує усталення; на скінченному записі перевіряють лише весь доступний подальший інтервал.
\end{frame}
\begin{frame}{Визначення: вихід 4--20 мА}
\[
 p=p_{\min}+\frac{I-4\;\mathrm{мА}}{16\;\mathrm{мА}}(p_{\max}-p_{\min}),
\]
\[
 I=4\;\mathrm{мА}+16\;\mathrm{мА}\frac{p-p_{\min}}{p_{\max}-p_{\min}}.
\]
За шкали 0--10 бар: 4 мА дає 0 бар, 12 мА дає 5 бар, 20 мА дає 10 бар. Зворотна підстановка відновлює струм.

Після переналаштування на 0--4 бар струм 12 мА означає 2 бар. Фізична здатність сенсора та налаштований розмах виходу є різними характеристиками.

Діагностичні струми й межі насичення визначає паспорт конкретного виконання. Не кожне число поза 4--20 мА має однаковий зміст.
\end{frame}
\begin{frame}{Промислові виконання й межі висновку}
WIKA S-20: сімейство промислових перетворювачів із різними діапазонами, виходами та приєднаннями.

Rosemount 3051: сімейство перетворювачів тиску та перепаду, конфігурації для рівня й витрати, діагностичні функції залежно від виконання.

Для вибору перевіряють: вид тиску, діапазон і перевантаження, матеріали змочуваних частин, ущільнення, температуру, швидкодію, живлення, вихід, приєднання й потрібне захисне виконання.

Розділювач і капіляр відокремлюють сенсор від середовища, але додають власні температурні й динамічні впливи. Числа навчальних паспортів наприкінці не є специфікаціями цих виробників.
\end{frame}
\begin{frame}{Історії та перевірені джерела}
\scriptsize
\begin{itemize}
\item Buncefield: Why did it happen? Competent Authority, 2011.\newline\url{https://www.icheme.org/media/10706/buncefield-report.pdf}
\item USGS. How Streamflow is Measured.\newline\url{https://www.usgs.gov/water-science-school/science/how-streamflow-measured}
\item WIKA. Bourdon tube pressure gauge.\newline\url{https://www.wika.com/en-en/lp_bourdon_tube_pressure_gauge.WIKA}
\item NASA. Experiencing a Dust Devil, 13.12.2022.\newline\url{https://science.nasa.gov/blog/experiencing-a-dust-devil/}
\item Endress+Hauser. Radar level measurement explained.\newline\url{https://www.endress.com/en/support-overview/learning-center/radar-level-measurement-explained}
\end{itemize}
Перевірено 25.10.2026. Фото, події та математичні навчальні моделі мають різне походження.
\end{frame}
\begin{frame}{Технічні джерела та екзаменаційне покриття}
\tiny
\url{https://www.wika.com/en-us/s_20.WIKA}\par
\url{https://www.emerson.com/en/measurement-instrumentation/products/rosemount-3051-pressure-transmitter-family}\par
\url{https://www.kistler.com/INT/en/piezoelectric-pressure-sensors-with-charge-output-pe-for-static-pressure-measurement/C00000140}
\small
\begin{enumerate}\setcounter{enumi}{51}
\item Механічні та електромеханічні методи вимірювання рівня рідини.
\item Радарні та гідростатичні рівноміри.
\item Рідинні, механічні деформаційні, електромеханічні та електрофізичні манометри.
\item Інтелектуальні вимірювальні перетворювачі тиску.
\item Промислові типи вимірювальних перетворювачів тиску.
\end{enumerate}
Формулювання за офіційним переліком 2020 року; питання 54 охоплюють блоки 3 і 4 разом.
\end{frame}
\appendix
\begin{frame}{Практичне завдання 1. Поплавковий канал}
\small
Напишіть функцію, що за струмом лінійного поплавкового каналу повертає рівень і статус діапазону. Синтетична калібровка: 4--20 мА відповідає 0--2 м.
\begin{itemize}
\item Вхідний список: \texttt{[3.5, 4, 8, 12, 20, 21, None]}; одиниця числових показів --- мА.
\item Для числового показу збережіть обчислений рівень у метрах та статус; для пропущеного поверніть відсутнє значення і окрему причину.
\item Межі включні. Не обрізайте покази мовчки й не змішуйте вихід за шкалу з відсутністю даних.
\item Перевірте обидві крайні точки, середину та обернене перетворення. Укажіть допуск порівняння чисел із рухомою комою.
\item Поясніть, чому перевірка діапазону не виявляє застрягання механізму на допустимому значенні.
\end{itemize}
Засоби Python: функція, список, словник, умовні оператори. Рішення та псевдокод не надаються.
\end{frame}
\begin{frame}{Практичне завдання 2. Порівняння рівнемірів}
\small
Поверніть $h_p$, $h_t$, $r=h_p-h_t$ і список попереджень для синтетичних даних: $\Delta p=19620$ Па, $\rho=1000$ кг/м$^3$, $g=9{,}81$ м/с$^2$, $H=3$ м, $v=343$ м/с, $t=0{,}005831$ с.
\begin{itemize}
\item Повторіть розрахунок із $\rho=800$ кг/м$^3$, решта вхідних даних незмінна; поясніть зміну лише одного каналу.
\item Додайте закритий бак: $p_{\mathrm{низ,g}}=49620$ Па, $p_{\mathrm{газ,g}}=30000$ Па; спершу визначте перепад.
\item Перевірте додатність $\rho,g,v$, невід'ємність часу та фізичні межі $0\le h\le H$; окремо обробіть відсутній опорний тиск.
\item Позначте $|r|>0{,}05$ м як навчальний привід для перевірки; не називайте його статистичним тестом.
\item Зворотно обчисліть тиск і час для граничних рівнів; поясніть роль мертвої зони.
\end{itemize}
Не усереднюйте оцінки автоматично. У звіті запропонуйте фізичні причини розбіжності.
\end{frame}
\begin{frame}{Практичне завдання 3. Шкали манометрів}
\small
Для робочого тиску 2,5 бар і вимоги $B\le0{,}07$ бар перевірте три синтетичні паспорти: 0--4, 0--6, 0--10 бар, усі з класом 1,5 за правилом відсотка від розмаху.
\begin{itemize}
\item Представте кандидатів словниками з ID, нижньою й верхньою межами та класом.
\item Поверніть для кожного розмах, абсолютну межу, результат перевірки діапазону й точності та текстову причину відхилення.
\item Не округляйте перед перевіркою. Окремо протестуйте рівність межі вимозі та рівність тиску кінцю шкали.
\item Додайте некоректну шкалу з однаковими кінцями, від'ємний клас і тиск поза діапазоном.
\item У висновку перерахуйте характеристики, яких бракує для реального вибору: середовище, температура, пульсації, перевантаження та приєднання.
\end{itemize}
Належність діапазону сама по собі не доводить повної придатності.
\end{frame}
\begin{frame}{Практичне завдання 4. Відбір перетворювачів}
\scriptsize
Вимога: весь діапазон 0--6 бар, температура 60~°C, вода, час відгуку $t_{90}\le0{,}10$ с. Усі наведені паспорти синтетичні, час виміряно за однаковою процедурою.
\begin{center}\begin{tabular}{ccccc}\toprule
ID & Тиск, бар & Температура, °C & Середовище & $t_{90}$, с \\\midrule
A & 0--10 & $-20$--80 & вода & 0,05 \\
B & 0--4 & $-20$--80 & вода & 0,03 \\
C & 0--10 & $-20$--50 & вода & 0,05 \\
D & 0--10 & $-20$--80 & невідомо & 0,04 \\
E & 0--10 & $-20$--80 & вода & 0,20 \\\bottomrule
\end{tabular}\end{center}
\small
Створіть функцію відбору зі статусами «підходить», «не підходить», «недостатньо даних». Збережіть результати кожного критерію та пояснення. Явна відмова має пріоритет над невідомим полем, але невідомі поля також перелічіть.

Перевірте граничний час 0,10 с, відсутній параметр, сумісність середовища й повне включення потрібного діапазону. Поясніть, чому $t_{90}$ не можна без перерахунку порівнювати з часом усталення.
\end{frame}
\end{document}
